% !TEX program = xelatex
% Generic Huawei Cup modeling manuscript starter.
% Formatting: automatically check the CURRENT official notice/template at WRITE time.
% Enforce only explicit mandatory requirements; treat recommendations as warnings.
% Do not infer unspecified rules from old years or from visual appearance alone.
% If the user/official source supplies a LaTeX template/class, preserve its explicit required structure.
\documentclass[UTF8]{ctexart}
\usepackage{amsmath,amssymb,amsthm,booktabs,graphicx,float}
\usepackage[hidelinks]{hyperref}

% Neutral working default only; user instructions override it.
\pagestyle{plain}

\begin{document}

\begin{center}
  {\Large 论文题目}
\end{center}

\noindent\textbf{摘要：}用简洁文字说明问题、模型、方法、主要结果、验证与结论边界。摘要的具体位置、篇幅、字号等仅按当届官方明确 mandatory 或用户明确要求设置。

\noindent\textbf{关键词：}关键词一；关键词二；关键词三

\section{问题重述}
只保留理解任务所需的题意、对象和目标。

\section{问题分析与建模思路}
说明题目结构为什么导向当前表示和模型；若确有实质替代且比较有助于论证，再简述选择依据，不为形式制造替代路线。

\section{模型假设与符号说明}
给出有依据的假设和统一符号。

\section{模型建立与求解}
集中给出变量、目标、约束、关键公式与求解方法，并解释方法为何适合模型。

\section{结果与验证}
给出回答题目的关键结果，并紧邻展示误差、基线、约束检查、敏感性或其它题型匹配的验证证据。

\section{结论、局限与推广}
把结果翻译回题目语言，写清支持什么、不支持什么以及适用范围。

\begin{thebibliography}{9}
\bibitem{source} 实际使用并核验过的可靠来源。
\end{thebibliography}

\end{document}
